\section{Fixed targets, factors and exchanges}\label{sec:targets}

This section gives applications of classical matching and factor
theory. These applications distinguish three tasks: independently
selecting mutually distinguishing targets, realising one target vector
by shared edge weights, and finding a target for which the actual
colour constraints hold. Only the last task for every graph would
prove the general conjecture.

\subsection{Independent target lists are not edge realisations}

\begin{lemma}\label{lem:complementlists}
Give each vertex $v$ of a simple graph $H$ a real number $D_v$ and a
list $L_v$ of at least $d_H(v)+1$ distinct real numbers. One can
choose $z_v\in L_v$ such that both $z_v$ and $D_v-z_v$ distinguish
the endpoints of every edge.
\end{lemma}
\begin{proof}
Orient the edges arbitrarily and write
$Q(z)=\prod_{uv}(z_u-z_v)=\sum_\alpha q_\alpha z^\alpha$.
For $d=(d_H(v))_v$ and $m=|E(H)|$,
$z^d Q(z^{-1})=(-1)^m Q(z)$, so
$q_{d-\alpha}=(-1)^m q_\alpha$. The coefficient of $z^d$ in
the highest homogeneous part of $Q(z)Q(D-z)$ is therefore
$\sum_\alpha q_\alpha^2>0$. Polynomial interpolation on the lists,
equivalently the Combinatorial Nullstellensatz~\cite{Alon99}, gives
a nonzero evaluation.
\end{proof}

This lemma says nothing about incidence realisation. On a triangle,
the lists $\{0,2,4\}$ with $D_v=4$ allow mutually distinct
complementary targets, but the targets $0,2,4$ cannot be realised
by weights $0,1,2$: the zero vertex forces its two weights zero,
preventing its degree-four neighbour from receiving four.

\subsection{A precise fixed-target criterion}

Consider a simple subgraph $H$ whose vertices receive fixed external
counts $a(v),b(v)$. Specify total red targets $R(v)$ and put
\[
 q(v)=R(v)-a(v),\qquad B(v)=b(v)+2d_H(v)-q(v).
\]
Define $E_0=\{uv:R(u)=R(v)\}$ and
$E_2=\{uv:B(u)=B(v)\}$. If the two sets overlap, no legal realisation
of these targets exists. Otherwise set
\[
 K=H-(E_0\cup E_2),\qquad t(v)=q(v)-2d_{E_2}(v).
\]

\begin{proposition}\label{prop:fixedfactor}
The targets admit a legal internal colouring exactly when
$q(v)\in[0,2d_H(v)]\cap\mathbb Z$, $E_0\cap E_2=\varnothing$,
and $\doub K$ has a spanning subgraph of degrees $t(v)$.
\end{proposition}
\begin{proof}
On $E_0$ all weights are forced to zero, on $E_2$ to 2. Every
remaining edge distinguishes both total colour degrees and is free
to have any weight. Its selected red copies must supply exactly $t$.
This correspondence is reversible.
\end{proof}

Here is an explicit ordinary matching reduction. For each copy of
$\doub K$ create its two endpoint ports and join them. At $v$ create
$2d_K(v)-t(v)$ auxiliary vertices adjacent to all ports at $v$.
A perfect matching uses that many local ports, leaving exactly
$t(v)$ ports matched across original copies. Select those copies
for the factor. Conversely any such factor completes to a perfect
matching by matching unused ports locally. The auxiliary graph has
at most $8|E(K)|$ vertices. Range failures for $t$ are rejected first.

\subsection{Capacity and odd-component certificates}

For disjoint $S,T\subseteq V(K)$ put
\[
 g_t(S,T)=t(S)+2\sum_{v\in T}d_{K-S}(v)-t(T),
\]
and let $o_t(S,T)$ be the number of components $C$ of $K-S-T$ for
which $t(C)$ is odd.

\begin{proposition}\label{prop:capacity}
Real weights $0\le x_e\le2$ realising $t$ exist if and only if
$g_t(S,T)\ge0$ for all disjoint $S,T$. Integer weights exist if and
only if $g_t(S,T)\ge o_t(S,T)$ for all such pairs.
\end{proposition}
\begin{proof}
Necessity of capacity follows by subtracting the sums over $T$ and
$S$ and bounding the remaining terms by
$4|E_K(T)|+2|E_K(T,V\setminus(S\cup T))|$.
For sufficiency, let $\mathcal Z=\{Bx:x\in[0,2]^{E(K)}\}$, where
$B$ is the unsigned incidence matrix. Its support function at $y$ is
$2\sum_{uv}\max(0,y_u+y_v)$. If $t\notin\mathcal Z$, a separating
vector $y$ violates this bound. Write $y$ as the integral, over
$h\ge0$, of the vector taking value 1 where $y_v>h$, value $-1$
where $y_v<-h$, and zero elsewhere. For each edge the same integral
identity holds for the positive part of $y_u+y_v$; it is checked
separately when the two signs agree or disagree. Some three-valued
vector must therefore violate the support bound, which is precisely
a violated $g_t(S,T)\ge0$.

For integers apply Tutte's factor theorem~\cite{Tutte54} to $\doub K$.
In its odd-component term the number of edges from a component to
$T$ is even because every original edge is doubled. The odd term
is consequently exactly $o_t(S,T)$, and its degree term is $g_t$.
This proves the stated equivalence.
\end{proof}

For even total target sum, $g_t-o_t$ is even. Colour complementation
sends $t$ to $2d_K-t$ and interchanges $S,T$, preserving both
quantities. If every $t(v)$ is even, every odd-component term is
zero, so continuous feasibility already gives an integer realisation.
One can also see this directly from a vertex of the capacity polytope.
Strictly fractional edges have independent incidence columns, hence
each of their components is a tree or has one odd cycle. Leaves
are impossible with even residual demands and all other weights at
0 or 2. The fractional support is thus a disjoint union of odd cycles;
each adjacent pair sums to 2, forcing all their weights to 1.
The resulting extreme point is integer. This is a fixed even-target
integrality fact, not permission to require all red degrees even in
the original conjecture.

Prescribed parity intervals also reduce to matching. Tighten
$[\ell(v),u(v)]\cap[0,2d(v)]$ to its first and last values
$g(v),f(v)$ of the desired parity. Add $(f(v)-g(v))/2$ virtual loops
at $v$ to $\doub G$. An exact factor of degree $f$ becomes the
required interval assignment after virtual loops are deleted, and
every such assignment can be completed by loops. For a matching
gadget with loops, create two ports per edge or loop, join each port
pair, and create $f(v)$ demand vertices joined to all local ports.
Ports are either matched together (edge unselected) or both to
demands (selected); exactly $f(v)$ incidences are selected. This
polynomial reduction includes degree feasibility only.

\subsection{Minimum conflicts and fixed-degree exchanges}

Let $\mathcal D(G)$ be all realisable integer red target vectors, and
\[
 \Psi(x)=\sum_{uv}\left(x_{uv}\mathbb{1}_{R(u)=R(v)}
 +(2-x_{uv})\mathbb{1}_{B(u)=B(v)}\right),\qquad
 F(q)=\min_{Bx=q}\Psi(x).
\]
For a fixed $q$ put $I_R(uv)=\mathbb{1}_{q(u)=q(v)}$,
$I_B(uv)=\mathbb{1}_{2d(u)-q(u)=2d(v)-q(v)}$. Then
\begin{equation}\label{eq:fixedcost}
 F(q)=2\sum_e I_B(e)+\min_{Bx=q}\sum_e (I_R(e)-I_B(e))x_e.
\end{equation}
The matching gadget above, with weight $-(I_R-I_B)$ on each
cross-port pair and zero on local edges, computes the minimum by
a maximum-weight perfect matching. This is an exact application of
weighted matching; the costs change when $q$ changes.

If two assignments realise the same target, mark their differing
copies by $+$ or $-$ and pair opposite signs at each vertex. They
decompose into even alternating closed walks. Altering one walk
moves toward the other assignment, never leaves $\{0,1,2\}$, and
preserves all degrees. Repeated vertices and repeated original edges
must be allowed. At fixed targets a legal realisation, if it exists,
can be reached without increasing $\Psi$, since all forced edges
move monotonically to their required extremes and all free edges
permit both colours.

More generally, for $q,r\in\mathcal D(G)$ and any unit step $s$
toward $r$, there is a second unit step $t$ toward $r$ such that
\[
 q+s+t,\ r-s-t\in\mathcal D(G),\qquad
 \|q+s+t-r\|_1=\|q-r\|_1-2.
\]
To prove this, pair opposite marked incidences as far as possible.
Unpaired incidences have exactly the signs of $r(v)-q(v)$. Start at
one representing $s$, follow its alternating paired trail, and stop
at the next unpaired incidence, representing $t$. Transfer the marked
copies in opposite directions in the two assignments. Every edge
stays between its two initial values. This is the standard constant
parity jump-system exchange phenomenon; see~\cite{BC95}. In particular
the two-unit target graph is connected, with distance
$\|q-r\|_1/2$. For a \emph{fixed} linear edge-cost vector $c$ the
same transfer gives
\[
 f_c(q+s+t)+f_c(r-s-t)\le f_c(q)+f_c(r),\quad
 f_c(q)=\min_{Bx=q}c\cdot x,
\]
because the two assignments have conserved total cost. The function
$F$ in~\eqref{eq:fixedcost} does not share this conclusion in general;
Section~\ref{sec:targetbarriers} gives explicit failures.

\subsection{A parity-defect lower bound}

For $E_0\cap E_2=\varnothing$, contract every component of the free
graph $K$. Let $\mathcal O$ be the resulting vertices with odd $q$
sum. On a forced edge write $z_e=x_e$ for $e\in E_0$ and
$z_e=2-x_e$ for $e\in E_2$. Then $\Psi=\sum z_e$, and each free
component $C$ satisfies
\[
 q(C)\equiv\sum_{e\in\delta(C)}z_e\pmod2.
\]
After duplicate copies cancel modulo two, the bad forced copies
contain an $\mathcal O$-join in the quotient graph. Therefore
$F(q)$ is at least its minimum join cardinality. Pairing incidences
decomposes any join into paths between paired odd vertices and
closed parts; conversely shortest paths of a minimum perfect matching
on the odd vertices, cancelled modulo two, give a join. This proves
the familiar shortest-path matching formula for this lower bound.
It detects parity only; internal factor capacities can require a
larger defect.

\section{Normal forms and conditional structural reductions}\label{sec:reductions}

\subsection{Sparse mixed-edge normal form}

\begin{proposition}\label{prop:normalform}
Every assignment can be changed, preserving all colour degrees and
not increasing $\Psi$, until the weight-one edges form a cactus
whose cycles are vertex-disjoint odd cycles.
\end{proposition}
\begin{proof}
On an even closed trail of mixed edges alternately increase and
decrease weights by one. Both signs of this change are feasible and
preserve degrees. At fixed degrees $\Psi$ is linear, so one choice
does not increase it. Every traversed edge ceases to be mixed.
Repeat until no even closed trail remains. There is no even cycle.
Two cycles sharing at least two vertices give a theta subgraph, two
of whose three paths have the same parity, giving an even cycle.
Two odd cycles sharing exactly one vertex give an even closed trail.
Thus cycles are vertex-disjoint and odd, and all other edges of each
mixed component are bridges.
\end{proof}

If there are $n$ vertices, $c$ mixed components including isolates,
and $k$ mixed cycles, their edge count is $n-c+k\le n-c+\lfloor n/3\rfloor$.
This sparsification neither asserts nor needs $\Psi=0$.

\subsection{Energy cleanup and an exact one-edge identity}

Define $\mathcal E(x)=\sum_v(R(v)-d_G(v))^2$. If a present colour
has equal degree $k>\max(d(u),d(v))$ at both ends of $uv$, changing
one copy of that colour to the other changes the energy by
\[
 -4k+2(d(u)+d(v))+2\le-2.
\]
If a present colour has equal degree one, its edge is mixed. Changing
the other copy to that colour changes the energy by
$6-2(d(u)+d(v))\le0$, provided the graph has no isolated $K_2$
component. Both new endpoint colour degrees are even, so no new
degree-one conflict appears. Hence the lexicographic pair
$(\mathcal E,h_1)$, where $h_1$ counts bad degree-one copies, strictly
decreases under these two rules. Termination takes at most
$(\lfloor\mathcal E(x)/2\rfloor+1)(2m+1)-1$ steps. The remaining
bad colour degrees are between two and the larger endpoint original
degree. The mixed-edge normal form can subsequently be imposed,
preserving these statements because it fixes all degrees and changes
only edges that already had both colours.

The energy is not $\Psi$. On a four-vertex path, changing the word
$(0,1,0)$ to $(0,2,0)$ lowers $\mathcal E$ from two to zero while
increasing $\Psi$ from two to six. The all-mixed triangle has energy
zero and six bad copies and is already stopped by both rules.

For completeness there is an exact local accounting identity. For a
permitted change $x_e\mapsto x_e+\eta$, $e=uv$, $\eta=\pm1$, put
\begin{align*}
 H_z^\eta(e)=\sum_{zw\ne e}\bigl[&x_{zw}
 (\mathbb{1}_{R(w)=R(z)+\eta}-\mathbb{1}_{R(w)=R(z)})\\
 &+(2-x_{zw})(\mathbb{1}_{B(w)=B(z)-\eta}-\mathbb{1}_{B(w)=B(z)})\bigr].
\end{align*}
Then
\[
 \Psi(x')-\Psi(x)=\eta(\mathbb{1}_{R(u)=R(v)}-
 \mathbb{1}_{B(u)=B(v)})+H_u^\eta(e)+H_v^\eta(e).
\]
On $e$ the degree differences stay fixed and only multiplicities
change; the remaining affected edges are exactly those in the two
sums. Any claimed improvement must account for these neighbour terms.

\subsection{Exact elimination of peripheral trees}

If $G$ contains a cycle, its 2-core $C$ is nonempty and connected.
Every component outside it is a tree attached by exactly one edge:
a cycle outside or two attachments would enlarge the minimum-degree-two
subgraph. At a core vertex put $h=d_C(v)\ge2$ and let
$\ell,s,b$ count attached trees of orders one, two and at least three,
respectively. Let $k=\ell+s+b,D=h+k$. At a proposed final red degree
$q$ put $B=2D-q$ and define
\begin{align*}
 L_1(q)&=\{c\in\{0,1,2\}:c=0\Rightarrow B\ne2,
 \ c=1\Rightarrow q,B\ne1,\ c=2\Rightarrow q\ne2\},\\
 L_2(q)&=\{c\in\{0,1,2\}:c=0\Rightarrow B\ne4,
 \ c=2\Rightarrow q\ne4\},\\
 S_v(q)&=\ell L_1(q)+s L_2(q)+b\{0,1,2\},\qquad
 J_v(q)=S_v(q)\cap[q-2h,q].
\end{align*}
The sums are Minkowski sums, with $0L=\{0\}$. A graph colouring is
equivalent to core weights $y_e$ and targets $q_v$ satisfying
\[
 q_v-\sum_{e\in C,e\ni v}y_e\in J_v(q_v)
\]
and the actual red/blue comparison on every core edge with total
degree $D_v$. Necessity is restriction. For sufficiency split the
chosen sum into allowed bridge weights and apply the tree interface
lemma independently. Physical counts hold because
$q=t+\sum c_i\ge c_i$ and
$B=(2h-t)+\sum(2-c_i)\ge2-c_i$.

Each $S_v(q)$ is an integer interval or an interval of step two.
Indeed $L_1$ is an interval or $\{0,2\}$, and $L_2$ is an interval,
possibly $\{1\}$. A width of at least one in any other interval
fills the gaps of a sum of $\{0,2\}$ sets. This is an exact compact
description, not an approximation of feasible root degrees. Trees
of order at least three may therefore be replaced by any rooted
three-vertex path, preserving colourability in both directions.

\subsection{A sharp tree capacity extension}\label{sec:treecapacity}

\begin{theorem}\label{thm:treecapacity}
At a core vertex of core degree $h\ge2$, $k\ge h$ attached rooted
trees can be restored to a given legal colouring of the graph with
those trees removed, preserving every old edge colour.
\end{theorem}
\begin{proof}
Fix its old internal red degree $t\in[0,2h]$. The feasible final
targets are $\mathcal Q(t)=\{q\in[t,t+2k]:q-t\in J_v(q)\}$.
For $k\ge3$ this is the whole interval, except for: value 1 if all
branches are singletons and $t=0$; its colour-complement at $t=2h$;
value 4 if there are exactly three branches, all of order two, and
$t=0$; its complement at $t=2h$. To check this list only the low
targets $1,2,4$ and their complements need inspection. At $q=1$
only singleton bridge sets become $\{0,2\}$, causing a gap exactly
when all branches are singletons and $t=0$. At $q=2$ their total
upper bound is $2k-\ell\ge k\ge3$, sufficient for any $2-t$.
At $q=4$ the upper bound is $2k-s$, insufficient only when
$k=s=3,t=0$. All other targets allow all weights. Low and high
cases are disjoint since $h+k\ge5$.

If $j$ old incident edges are mixed, old neighbours forbid at most
$h+j$ new targets. When $0<t<2h$, there are $2k+1>h+j$ candidates.
At $t=0$ or $2h$ no old edge is mixed and at least $2k>h$ remain.
The small case $h=k=2$ has at least four candidates, and all five
at $t=2$: directly inspect $q=1,2,4,6,7$ in the bridge sets.
At $t=1,3$ exactly one old edge is mixed, giving at most three
forbidden targets; at $t=0,4$ there are at most two. Thus an allowed
target always exists. Split its bridge contribution and fill the trees.
\end{proof}

The bound cannot uniformly be reduced to $k=h-1$ for any $h\ge2$.
Create roots $v,u_1,\ldots,u_h,z$, edges $vu_i$, and, for each
$i$, $d_i-1$ internally disjoint two-edge paths from $u_i$ to $z$,
where $d_i=2h-2+i$. The old degrees are $h,d_i,5h(h-1)/2$ at
$v,u_i,z$ and two at the path interiors. All-mixed old edges are
legal. Attach $h-1$ three-vertex rooted paths at $v$. A preserved
old colouring has $q\in[h,3h-2]$ and total colour sum $4h-2$ there.
The old edges forbid $q=d_i$ and $q=4h-2-d_i$, covering that entire
interval. Yet a full recolouring exists: change only $vu_2$ to zero,
give new bridges weight 1 and each new internal path word $(0,0)$.
Now $v$ has $(2h-2,2h)$, $u_2$ has $(2h-1,2h+1)$, other
$u_i$ keep $(d_i,d_i)$, and each tree has successive counts
$(1,3),(0,4),(0,2)$. These values verify every edge.

\subsection{Universal vertices and minimality}

\begin{theorem}\label{thm:universal}
If $z$ is universal in $G\ne K_2$, and every component of $G-z$
of order at least three has a given legal colouring, then the
colouring extends to $G$, preserving those given edge weights.
Isolated vertices and $K_2$ components require no given colouring.
\end{theorem}
\begin{proof}
Let $H=G-z$ have $N$ vertices, $k$ components $K_2$, $\ell$
isolates and larger components of orders $c_i$. Join every larger
component and isolate to $z$ with blue edges. For each $K_2=uv$
put $uv=1,zu=0,zv=2$. Its endpoints become $(1,3),(3,1)$ and
contribute $(2,2)$ to $z$. Thus $z$ has
$(2k,2(N-k))$. The $K_2$ interfaces compare even counts to three;
an isolated blue neighbour has count two, equal only in the excluded
case $G=K_2$. In a larger component, all old blue degrees increase
by two, preserving every old comparison, and are at most $2c_i$.
If $H$ has any other component, $2(N-k)>2c_i$, so the new
interfaces are legal.

It remains that $H$ is a single connected component of order $N\ge3$.
All-blue new edges fail only if an old vertex is universal in $H$
and has every old copy blue. If no such vertex exists, retain them.
Otherwise make every new edge red. This preserves old differences
and can fail only at an old all-red universal vertex. An all-blue
and an all-red universal vertex cannot coexist, since their common
positive-multiplicity edge would need both pure colours. Hence the
second choice works.
\end{proof}

Consequently a vertex-minimal counterexample has no universal vertex,
so $\Delta\le n-2$. It has a nonempty 2-core, every attached tree
has at most three vertices, and the number $k_v$ of trees at a
core vertex is smaller than its core degree $h_v$. These follow by
replacement and Theorem~\ref{thm:treecapacity}; the smaller remainders
contain a cycle and are not $K_2$. Hence, for core order and size
$n_C,m_C$,
\[
 |V(G)|\le n_C+3\sum_v(h_v-1)=6m_C-2n_C.
\]
Every maximal degree-two core thread has length at most five,
including return threads. At length at least six replace its interior
and attached trees by a new bare two-edge path, or a triangle for a
return. The smaller graph is simple and connected and has a legal
colouring by minimality. Its fixed endpoint counts and end weights
extend back by Theorem~\ref{thm:decorated}. The condition $k_v<h_v$
ensures that each internal core-degree-two vertex has at most one
tree. If the core is a cycle, its length is also at most five.
For a noncycle core compressed to a skeleton with $n_H$ vertices,
$m_H$ edges and minimum degree at least three, there are at most
$4m_H$ degree-two vertices. Combining the branch bound gives
\[
 |V(G)|\le22m_H-2n_H.
\]
The skeleton size and degrees remain unbounded; none of these
necessary conditions resolves an arbitrary minimum-degree-three core.
